(1) 一个方向能量流为 \(J_+(T)=\int_0^\infty \hbar\omega n_B(\omega,T)\,d\omega/(2\pi)\)。换元 \(x=\hbar\omega/(k_BT)\)，利用
\(\int_0^\infty x/(e^x-1)\,dx=\pi^2/6\)，得到 \(J_+(T)=\pi k_B^2T^2/(12\hbar)\)。左右相减，
\[
\boxed{J_Q=\frac{\pi k_B^2}{12\hbar}(T_L^2-T_R^2)} .
\]
模态的 \(\hbar\omega/2\) 真空零点项对两个传播方向相同，且与热库温度无关，净输运中严格抵消，不能当作连续输出的热能。

(2) 对小温差展开 \(T_L^2-T_R^2=2T\Delta T+O(\Delta T^3)\)，所以
\[
\boxed{G_Q=\left.\frac{J_Q}{\Delta T}\right|_{\Delta T\to0}
=\frac{\pi k_B^2T}{6\hbar}=\frac{\pi^2k_B^2T}{3h}}.
\]
在色散严格线性且一维单分支无散射时，群速 \(v\) 与每单位长度的态密度 \(1/v\) 抵消，故结果不含声速。\(N\) 条无耦合完全透射分支给 \(NG_Q\)；它不是任何有限温度下三维扩散材料的普遍导热系数。

(3) 记 \(X_c=\hbar\Omega_c/(k_BT)\)，单向有限带宽流为
\[
J_+^{(c)}(T)=\frac{k_B^2T^2}{2\pi\hbar}
\int_0^{X_c}\frac{x\,dx}{e^x-1}.
\]
反导函数可取 \(x\ln(1-e^{-x})-\operatorname{Li}_2(e^{-x})\)；其在0的极限为 \(-\pi^2/6\)，所以积分
\(\pi^2/6+X_c\ln(1-e^{-X_c})-\operatorname{Li}_2(e^{-X_c})\)。
令 \(\Omega_c\to\infty\) 还原第一问。若 \(k_BT\gg\hbar\Omega_c\)，则 \(x/(e^x-1)=1-x/2+O(x^2)\)，积分给 \(J_+^{(c)}\simeq k_BT\Omega_c/(2\pi)\)，净流线性于 \(T_L-T_R\)。

(4) 若 \(T_L>T_R\)，定常无耗散通道中从热库L取热 \(J_Q\)，向R放热 \(J_Q\)。两个维持恒温的热库熵变化率和为
\[
\boxed{\dot S_{\rm tot}=J_Q\left(\frac1{T_R}-\frac1{T_L}\right)
=\frac{\pi k_B^2}{12\hbar}\frac{(T_L-T_R)^2(T_L+T_R)}{T_LT_R}\ge0}.
\]
式子在等温时为零，温差交换满足第二定律。有限传输概率 \(\mathcal T(\omega)\) 应放入积分核且 \(0\le\mathcal T\le1\)，此时不再自动具有独立于频谱的普适热导量子。
